\documentclass[10pt,a4paper]{article}
\usepackage[margin=20mm]{geometry}
\usepackage[utf8]{inputenc}\usepackage[T1]{fontenc}
\usepackage[dutch]{babel}\usepackage{lmodern}
\usepackage{amsmath,amssymb,booktabs,longtable,array,graphicx,xcolor,hyperref,fancyhdr,lastpage}
\hypersetup{hidelinks}\definecolor{main}{HTML}{173249}\definecolor{sub}{HTML}{176A75}
\newcommand{\rhoF}{\rho_{\!f}}\newcommand{\rc}{r_c}\newcommand{\vswirl}{\mathbf{v}_{\!\boldsymbol{\circlearrowleft}}}
\pagestyle{fancy}\fancyhf{}\fancyhead[L]{\small\color{sub}A075-v0.2.5 | Falsifier Framework v1.1.0}\fancyhead[R]{\small 28 oktober 2026}
\fancyfoot[C]{\small Pagina \thepage\ van \pageref{LastPage}}\setlength{\headheight}{15pt}
\title{\textcolor{main}{A075-v0.2.5}\\Hydrodynamische energie-invanging\\\large Driedimensionale Euler-dynamica, knoopgeometrie en operationele klok}
\author{SST-Workbench | Framework v1.1.0 adaptor (lokale runstatus in framework-ledger)}
\date{28 oktober 2026 | \textbf{DEVELOPMENT ONLY}}
\begin{document}\maketitle
\begin{abstract}
De onderzoeksvraag is of een vrij evoluerende, volledig gekoppelde wervelstructuur na verstrooiing extra lokale energie kan behouden gedurende veel eigen klokcycli, terwijl de lokale fluxbalans en numerieke convergentie sluiten. Deze release introduceert een zelfstandige, driedimensionale, periodieke Fourier--Galerkin-Eulerreferentie, met een finite-core-vorticiteitsinitialisatie, onafhankelijke ontvanger-zonder-broncontrole en twee onafhankelijke resolutierichtingen. Er zijn geen gefitte snelheids-, resonantie- of dempingsparameters. De uitgevoerde \texttt{PREVIEW} is bewust ondergeresolveerd en toont geen bronvertrek of onafhankelijke klok. \textbf{Fysische conclusie: UNRESOLVED.} De originele preview is zonder lokale Framework-FREEZE uitgevoerd. De latere lokale FREEZE- en gate-uitkomsten staan uitsluitend in de automatische framework-ledger.
\end{abstract}
\section{Exacte vraag en falsifieerbare nulhypothese}
\begin{quote}\emph{Kan een vrij evoluerende, volledig gekoppelde wervelstructuur na verstrooiing een positieve lokale energietoename behouden gedurende vele eigen klokcycli, terwijl de energiefluxbalans sluit en het resultaat numeriek convergeert?}\end{quote}
\textbf{$H_0$:} zolang bronvertrek, lange klokduur en onafhankelijke ruimtelijke, temporele en kern-convergentie ontbreken, is invanging niet aangetoond.\\
\textbf{$H_1$:} een uit onafhankelijke geometrie en vooraf gekozen parameters verkregen, positieve minimumretentie na geverifieerd vertrek over minstens 100 robuuste ontvangerperioden, met sluitende energieflux in een voldoende opgelost Euler-volume.
De eis van honderd cycli is hier een \emph{vooraf ingestelde operationele verificatiegrens}, geen nieuwe natuurconstante.
\section{Volledige veldvergelijkingen en druk}
Met constante dichtheid $\rhoF$, absolute evolutietijd $t$ en vorticiteit $\boldsymbol\omega=\nabla\times\mathbf{u}$:
\begin{align}
\nabla\cdot\mathbf{u}&=0, &
\partial_t\mathbf{u}+\mathbf{u}\cdot\nabla\mathbf{u}&=-\nabla p/\rhoF,\\
\partial_t\mathbf{u}&=\mathbb{P}(\mathbf{u}\times\boldsymbol\omega), &
\Delta\left(\frac p{\rhoF}+\frac{|\mathbf{u}|^2}2\right)&=\nabla\cdot(\mathbf{u}\times\boldsymbol\omega).
\end{align}
\noindent$\mathbb{P}$ is de divergente-vrije Lerayprojector. De initi\"ele drie gesloten kernen worden gedeponeerd als vorticiteitslijnen en door een expliciet vermelde Gauss-kern gladgemaakt. Het daaropvolgende 3D-veld volgt Fourier--Galerkin-Euler met tweederde-dealiasing en vierde-orde Runge--Kutta. \textbf{Geen cirkelgeometrie wordt tijdens evolutie opgelegd.} Deze numerieke benadering is geen bewijs van voortbestaan van drie afzonderlijke materiële tubes bij eindige roosterbreedte. De periodieke rand zorgt voor beeldinteracties en vereist een afzonderlijke domeinlengtestudie.
\section{Knoopconfiguraties en invoerherkomst}
\begin{center}\begin{tabular}{lll}\toprule
Test & Receiver 2 & Geometrische autoriteit\\\midrule
$0_1$ & circulaire unknot & analytische testgeometrie\\
$3_1$ & torustrefoil $T(2,3)$ & analytische embedding; niet ideaal\\
$5_2$ & twistknoop & alleen originele \texttt{KNOT\_5\_2\_N2048\_ORIG.xyz}\\
$6_1$ & twistknoop & alleen originele \texttt{KNOT\_6\_1\_N2048\_ORIG.xyz}\\\bottomrule
\end{tabular}\end{center}
In elk geval blijven de eerste ontvanger $0_1$ en de inkomende bron $0_1$. Alle drie hebben eindige, initieel gladde vorticiteitsbreedte. De twee externe twistknopen ontbreken in de geleverde omgeving en worden daarom niet gesimuleerd, noch vervangen door verzonnen knopen. Bij geladen invoerbestanden worden de oorspronkelijke bytes met SHA-256 vastgelegd; de topologische identiteit vereist vervolgens een afzonderlijke verificatie.
\section{Lokale energie, flux en controle zonder donor}
De lokale energiedichtheid en het volledige Eulerfluxveld luiden
\begin{equation}
 e=\frac{\rhoF}{2}|\mathbf{u}|^2,\qquad
\partial_t e+\nabla\cdot[(e+p)\mathbf{u}]=0.
\end{equation}
Voor een \emph{meebewegend maar vormvast} zacht venster $W(\mathbf{x}-\mathbf{X}(t))$ met venstersnelheid $\mathbf w=\dot{\mathbf X}$:
\begin{equation}
\boxed{\frac{d}{dt}\int W e\,dV=\int\nabla W\cdot\bigl[(e+p)\mathbf{u}-e\mathbf{w}\bigr]\,dV.}
\end{equation}
Het venster volgt een passieve referentiemarker van de \emph{controle zonder donor}; exact hetzelfde venster wordt voor de run met bron gebruikt. Dit voorkomt dat verschillende vensterposities een schijnbare verandering maken. \textbf{Het venster is geen volledig vervormend materieel volume.} Flux wordt vergeleken met $dE_W/dt$ en bij toekomstige kwalificatie verfijnd in $\Delta x$, $\Delta t$ en de vensterrandbreedte.
\begin{equation}
\Delta E_W(t)=E_W^{(\mathrm{bron})}(t)-E_W^{(\mathrm{controle})}(t),\qquad
\eta_{\mathrm{floor}}=\frac{\inf_{t\in W_{\mathrm{laat}}}\Delta E_W(t)}{E_{\mathrm{donor,isolated}}}.
\end{equation}
\textbf{Een positieve waarde bij begin of binnen een kort venster is geen absorptie}: de inkomende kern kan nog binnen het venster zijn, de ontvangende geometrie kan zijn verplaatst en de energie omvat interferentiekruistermen. $\eta_{\mathrm{floor}}$ is daarom pas een \emph{retentiemaat} na apart bronvertrek, latetime-klok- en fluxgate.
\section{Absolute tijd en ontvanger-eigenklok}
Een klok wordt niet vooraf ge\"identificeerd met de absolute evolutietijd. Eerst moet een onafhankelijk, stabiel periodiek ontvangstsignaal $q_A(t)$ worden gemeten. Zijn uitgepakte fase $\theta_A(t)$ geeft
\begin{equation}
 N_A[t_1,t_2]=\frac{\theta_A(t_2)-\theta_A(t_1)}{2\pi},\qquad
 \Delta\tau_A=\int_{t_1}^{t_2}\frac{d\tau_A}{dt}\,dt.
\end{equation}
Er is \emph{geen} automatische klokvalidatie op basis van enkele maxima. Deze release vereist bronvertrek, meer dan honderd gekwalificeerde ontvangercycli en robuuste piekperioden onder resolutieverandering. Een optisch-fotonmodel, atoomidentiteit of intrinsieke PKLSA-klok wordt niet aangenomen.
\section{Historische diagnose: preview uit de eerdere versie}
De onderstaande historische getallen komen uit de voorafgaande ontwikkelpreview, niet uit deze v0.2.5-release, \emph{niet} uit een authentieke Framework-FREEZE. De genormaliseerde eenheden zijn consequent interpreteerbaar als SI-schaalwaarden; er is geen atomair constante-fit gebruikt.
\begin{center}\begin{tabular}{lrr}\toprule
Diagnostiek & $0_1$ & $3_1$\\\midrule
Roostercellen per as (grof / fijner) & 24 / 32 & 24 / 32\\
Fysieke eindtijd (s) & 1.20 & 1.20\\
Kernbreedte per grof roosterpunt & 1.32 & 1.32\\
Maximale CFL-waarde & 0.07037 & 0.08188\\
Relatieve globale energiedrift & 1.955e-11 & 5.727e-11\\
Fluxbudgetresidu / beginenergie & 5.122e-05 & 5.857e-05\\
Ruimtelijke verschilmaat / donorenergie & 0.00389 & 0.01271\\
Temporele verschilmaat / donorenergie & 4.306e-06 & 9.908e-06\\
Gekwalificeerde klokperioden & geen & geen\\\bottomrule
\end{tabular}\end{center}
De geringe globale energiedrift is een noodzakelijke conserveringscontrole, maar de kernbreedte op het grove rooster is slechts 1.32 rasterstappen; \textbf{beide previewcases zijn dus ondergeresolveerd}. De lokale controleverandering is \emph{niet} een post-exit absorptiemeting. Bronvertrek werd niet gemeten. De onbekende fysische uitkomst is geen negatieve of positieve ontdekking.
\section{Wetenschappelijke gates, stopcriteria en framework-integratie}
\begin{longtable}{@{}p{10mm}p{25mm}p{111mm}@{}}\toprule
Gate & Status & Reden en minimale vervolgeis\\\midrule\endfirsthead\toprule Gate & Status & Reden\\\midrule\endhead
G0 & NOT RUN & Authentieke Framework v1.1.0 FREEZE is op dit systeem niet gedraaid. Nieuwe private commitment vereist.\\
G1 & PASS & A021/A032-methodologiebytes afzonderlijk tegen SHA-256 gecontroleerd.\\
G2 & UNRESOLVED & Tijd- en ruimterefinement aanwezig, maar kern onvoldoende opgelost; derde niveau en core- en domeinconvergentie ontbreken.\\
G3 & UNRESOLVED & Geen bewezen bronvertrek, geen robuuste eigenklok en geen honderd-cycli-retentie.\\
G4 & NOT RUN & Geen vooraf geregistreerde onafhankelijke blind-holdout.\\
G5/G6 & UNRESOLVED & De v1.1.0-backenduitvoering wordt in de runmatrix vastgelegd; dit bevestigt niet zelfstandig de fysische geldigheid.\\
G7 & BLOCKED & De originele $5_2$/$6_1$ geometriebestanden ontbreken; onafhankelijke topologiecontrole nodig.\\
G8 & UNRESOLVED & Periodiek Euler-veld bestaat, maar finite-core-, venster-, domein- en topologieconvergentie ontbreken.\\
G9 & WITHHELD & Geen fysieke statuspromotie of private reveal.\\\bottomrule\end{longtable}
\noindent\textbf{Exacte gebruikersvolgorde op Windows:} \texttt{python -m pytest -q tests}; \texttt{run\_all.cmd FREEZE}; \texttt{run\_all.cmd BASIC}; \texttt{run\_all.cmd FULL}. Met de ingebedde Framework-runtime is rechtstreeks \texttt{python run\_preview.py --profile PREVIEW} worden uitgevoerd. Die stand-alone run is geen Framework-gatevalidatie en kan geen nieuwe frozen protocol-IDs uitgeven. De originele A075-v0.1.0 bestanden en outputs blijven ongewijzigd.
\section{Eindantwoord}
\begin{quote}\textbf{UNRESOLVED.} De huidige data bewijzen geen langdurige invanging in een volledig convergente driedimensionale Eulerstroming. De nieuwe release biedt wel een werkelijke 3D-veldprovider, no-donorcontrole, druk-energieboekhouding, reproduceerbare knoopselectie en aparte temporele/ruimtelijke vergelijkingen. De niet-opgeloste bronvertrek-, klok-, kern-, domein- en holdoutgates verhinderen iedere positieve fysische claim.\end{quote}
\section{Automatisch vastgelegde registratie en gate-uitkomsten}
Onderstaande gegevens worden uitsluitend tijdens een echte Framework-uitvoering gegenereerd.
Bij afwezigheid van bestanden is geen framework-gate uitgevoerd.
\IfFileExists{AUTO_SCIENCE_CONTRACT.tex}{% AUTO-GENERATED. Do not edit; generated from frozen science_contract.json.
\subsection*{Machine-readable science contract}
\begin{description}
\item[Research question] Can a source-free, fully 3D Euler vorticity-field receiver, composed of closed finite-core rings/knots, retain additional local kinetic energy through an independently observed post-source-exit interval spanning many receiver clock cycles, while local energy flux and convergence close?
\item[Objective] Adversarial falsification: 3D periodic Euler screening with separate controls for conservative energy, pressure, comoving control volume, explicit 0\_1/3\_1/5\_2/6\_1 embeddings, and strong non-promotion of short run proxies.
\item[$H_0$] H0: no independent qualification of source departure, positive late minimum energy above no-donor control across many cycles and convergent energy-flux closure.
\item[$H_1$] H1: all preregistered post-exit late-time, core, spatial, time, periodic-box, clock and no-donor controls pass in an independent holdout.
\end{description}
\subsubsection*{Registered equations}
\paragraph{E1: Euler and incompressibility}
\begin{equation}
\partial_t\mathbf{u}=\mathbb{P}(\mathbf{u}\times\boldsymbol\omega),\ \nabla\cdot\mathbf{u}=0,\ \boldsymbol\omega=\nabla\times\mathbf{u}
\end{equation}
\noindent\textit{Dimensional check:} L/T\textasciicircum{}2\\
\paragraph{E2: Pressure-Poisson recovery}
\begin{equation}
\Delta(p/\rho+|\mathbf{u}|^2/2)=\nabla\cdot(\mathbf{u}\times\boldsymbol\omega)
\end{equation}
\noindent\textit{Dimensional check:} 1/T\textasciicircum{}2\\
\paragraph{E3: Local energy balance translating fixed-shape window}
\begin{equation}
\frac{d}{dt}\int W e\,dV=\int\nabla W\cdot[(e+p)\mathbf{u}-e\mathbf{w}]\,dV,\quad e=\rho|\mathbf{u}|^2/2
\end{equation}
\noindent\textit{Dimensional check:} J/s\\
\paragraph{E4: Matched donor versus no-donor windows}
\begin{equation}
\Delta E_W(t)=E_W^{\rm donor}(t)-E_W^{\rm control}(t)
\end{equation}
\noindent\textit{Dimensional check:} J\\
\paragraph{E5: Only post-exit retentive floor}
\begin{equation}
\eta_{\rm floor}=\inf_{t\in W_{\rm late}}\Delta E_W(t)/E_{\rm in}
\end{equation}
\noindent\textit{Dimensional check:} 1\\
\paragraph{E6: Separate absolute time and local receiver clock cycles}
\begin{equation}
N_A=\frac{\theta_A(t_2)-\theta_A(t_1)}{2\pi}
\end{equation}
\noindent\textit{Dimensional check:} 1\\
\subsubsection*{Registered code-to-mathematics steps}
\begin{enumerate}
\item \textbf{S01} [G0; equations E1]: Freeze authentic Framework v1.1.0 protocol locally before scientific runs
\item \textbf{S02} [G2; equations E1, E3, E4]: Run baseline without donor and matched donor run with identical precommitted control-window trajectory
\item \textbf{S03} [G2; equations E1, E3, E4]: Repeat at half timestep holding the absolute physical final time and geometry fixed
\item \textbf{S04} [G2; equations E1, E2, E3]: Repeat at finer grid spacing holding the physical core radius and final time fixed
\item \textbf{S05} [G3; equations E3, E5, E6]: Check periodic box, core profile, source exit by material transport, and independent clock qualification
\item \textbf{S06} [G4; equations E4, E5, E6]: Execute independent held-out initial conditions and source-qualified knot geometries after locked gates
\end{enumerate}
}{\par\textit{Science-contract: niet aanwezig in deze renderer.}}
\IfFileExists{AUTO_RESULTS.tex}{\input{AUTO_RESULTS.tex}}{\par\textit{Run-resultaten: niet aanwezig in deze renderer.}}

\section{Framework v1.1.0: wetenschappelijk volledigheidscontract}
% SST-REPORT-SECTION:IDENTIFICATION
\subsection{Framework- en pakketidentiteit}
De onderzoeksfamilie A075 wordt uitgevoerd als afzonderlijke versie v0.2.5. Het onderliggende Framework v1.1.0 is opgenomen als eigen, byte-gecontroleerde runtime-snapshot met vastgelegde bronhashes.
% SST-REPORT-SECTION:QUESTION
\subsection{Centrale falsifiervraag}
De primaire vraag is of werkelijk langdurige lokale kinetische energie kan achterblijven in de ontvanger na onafhankelijk vastgesteld bronvertrek, gedurende voldoende interne klokcycli en met gesloten energieflux.
% SST-REPORT-SECTION:HYPOTHESES
\subsection{Nul- en alternatieve hypothese}
De nulhypothese betreft het ontbreken van onafhankelijk gekwalificeerde retentie. Een alternatieve hypothese vergt post-exit meetvensters, vooraf vastgelegde drempels, convergentie en een onafhankelijke holdout.
% SST-REPORT-SECTION:ASSUMPTIONS
\subsection{Fysische en numerieke aannames}
De stroming is ideaal, inkompressibel en met constante dichtheid. De numerieke Fourier-Galerkin-representatie is periodiek; een vergroting van de rekenbox is geen vervanging voor een vrije-ruimte-oplosser.
% SST-REPORT-SECTION:SYMBOLS
\subsection{Gebruikte grootheden}
De snelheid heeft eenheid meter per seconde, vorticiteit inverse seconde en kinetische energie joule. De kernbreedte betreft een initiële Gaussische vorticiteitsschaal en is geen scherp gedefinieerde kernradius.
% SST-REPORT-SECTION:EQUATIONS
\subsection{Vergelijkingen en dimensionele controle}
De Eulervergelijking, de Fourier-Lerayprojectie, de drukrelatie en de bewegende-vensterflux worden in de voorafgaande secties gegeven. De dimensies van de energiebalans zijn overal joule per seconde.
% SST-REPORT-SECTION:ALGORITHM
\subsection{Reproduceerbaar algoritme}
Elke numerieke stap gebruikt de vierde-orde Runge-Kutta-integrator met adaptieve stapbeperking, eindpuntcontrole en afzonderlijke globale energiedrift. Bron- en controletraject delen dezelfde tijdroosters.
% SST-REPORT-SECTION:SOURCES
\subsection{Brongeometrie en herkomst}
De unknot en trefoil hebben expliciete analytische invoerdatasets. Voor de twistknopen vijf-twee en zes-een blijven authentieke externe XYZ-bestanden verplicht; ontbrekende bronbestanden worden niet vervangen.
% SST-REPORT-SECTION:GATES
\subsection{Gekwalificeerde en open gates}
Software-uitvoering, rekenprecisie en backendidentiteit worden afzonderlijk van fysische retentie geaudit. Geen positieve retentieclaim is toegestaan zolang bronvertrek, honderd klokcycli of onafhankelijke fluxconvergentie ontbreken.
% SST-REPORT-SECTION:NUMERICS
\subsection{Numerieke verfijningsstrategie}
Basale en verfijnde roosters worden bij dezelfde fysieke eindtijd gebruikt. De eindige-kernresolutie, tijdstap, doosgrootte en vensterrand moeten als afzonderlijke convergentieparameters worden gevarieerd.
% SST-REPORT-SECTION:BACKENDS
\subsection{Drie echte rekenroutes}
De Pythonroute gebruikt SciPy voor Fouriertransformaties, de C++-route voert de niet-lineaire kern via pybind11 uit, en de Intelroute voert de spectrale Eulerstappen in PyTorch op een werkelijk XPU-apparaat uit.
% SST-REPORT-SECTION:BLINDNESS
\subsection{Blind protocol en behoud}
De versie bevat geen eerder gebruikte private sleutel, onthullingsbestand of bevroren protocol. Het Framework mag alleen na volledige structurele preflight een nieuw blind protocol bevriezen; bestaande versies worden niet gewijzigd.
% SST-REPORT-SECTION:STATISTICS
\subsection{Meetbaarheden en fouten}
De diagnostische observabelen zijn eindige vensterenergie, minimale late energie en afzonderlijke globale energiedrift voor bron- en controlerun. Een verschil tussen backendobservabelen is nog geen aantoonbaar fysisch effect.
% SST-REPORT-SECTION:RESULTS
\subsection{Status van deze nieuwe versie}
De broncode en de frameworkconfiguratie worden in deze ontwikkeluitgave geleverd. Reële Windows-uitvoering met Intel Arc A770 en volledige convergentietoetsen zijn niet in de distributie uitgevoerd en mogen niet als PASS worden behandeld.
% SST-REPORT-SECTION:INTERPRETATION
\subsection{Wetenschappelijke interpretatie}
Een geslaagd twee-uit-drie-uitvoeringscontract toont dat twee identificeerbare wetenschappelijke providers werkelijk hebben gerekend. Het toont geen vrije-ruimte-limiet, kernconvergentie, fysieke klok of langdurige energie-invanging.
% SST-REPORT-SECTION:REPRODUCIBILITY
\subsection{Herhaalbaarheid en audittrail}
De versie legt de gebruikte softwarebronnen, runtime-snapshot, input-identiteit, uitvoerbare backend en werkelijk berekende meetvectoren met SHA-256 vast. Private onthullingsgegevens blijven buiten het blinde archief.

\begin{thebibliography}{9}
\bibitem{Saffman} P. G. Saffman, \emph{Vortex Dynamics}, Cambridge University Press (1992).
\bibitem{A021} A021-v0.4.0, Trefoil Lobe Self-Confinement, SST-Workbench methodedossier; bijgevoegd als bytes en SHA-256.
\bibitem{A032} A032-v0.1.0, Kelvin--Joule Transient Energy, SST-Workbench methodedossier; bijgevoegd als bytes en SHA-256.
\bibitem{A075} A075-v0.1.0 en afzonderlijke Hamiltoniaanse drieringspilot, verkennende modellen; geen complete Eulerkwalificatie.
\end{thebibliography}

\section{Adaptive Time-Stepping and Backend Evidence Matrix (v0.2.5)}
For each numerical time step $n$ the signed global energy drift is logged as
\begin{equation}
\delta E_n=E(t_{n+1})-E(t_n),\qquad
\epsilon_E(t_n)=\frac{E(t_n)-E(0)}{|E(0)|}.
\end{equation}
The numerical energy lost or gained per step is a \emph{numerical error diagnostic}:
$\max(0,-\delta E_n)$ and $\max(0,\delta E_n)$. In a closed, inviscid,
periodic Euler model these are not physical dissipation or absorption.
Given maximum speed $U_{\max}$, maximum vorticity magnitude $\omega_{\max}$,
mesh spacing $\Delta x$, Gaussian core scale $a$ and absolute time remaining,
\begin{equation}
\Delta t_n=\min\left(\Delta t_{\max},\,
C_{\mathrm{CFL}}\frac{\Delta x}{U_{\max}},\,
C_a\frac{a}{U_{\max}},\,
\frac{\theta_{\max}}{\omega_{\max}},\,T_f-t_n\right).
\end{equation}
Here the vorticity bound is a rotation \emph{proxy}, not a full strain bound;
a separate convergence study is essential. The final step is clipped to $T_f$,
and donor/control states use an identical adaptive schedule.
The source package provides per-step CSV energy matrices at base,
half maximum time step, and refined grid. Comparison across Python/SciPy,
C++/pybind11/OpenMP (nonlinear kernel; SciPy FFT), and optional real
PyTorch XPU (device-resident Fourier, projection, pressure, and RK4) is
possible only when each provider actually executes. XPU unavailable or
FP32-only parity cannot be promoted to certified native FP64 parity.
\textbf{A075 scientific conclusion remains UNRESOLVED} until donor departure,
receiver clock, local flux, finite-core, box-size and time/space convergence
are all independently demonstrated.

\end{document}
